% SGA 6, Exposición VII: fragmento español alineado con la fuente.
% Alcance: apertura y sección 1 completas, líneas 9752--9964 inclusive.
% Fuente francesa: sga6_fr_r1/src/fr.tex.
% SHA-256 de la fuente francesa completa:
% 0D2158B9FE3B466309A92165727A0B191422E05FBAA0DACD5806E8322C54BA1A.
\providecommand{\sisi}[2]{#2}

\section*{Exposición VII. Inmersiones regulares y cálculo del \texorpdfstring{$K^{\bullet}$}{K0} de un esquema explotado}
\begin{center}
por Pierre \textsc{Berthelot}
\end{center}

\section*{1. Generalidades sobre las inmersiones regulares}

Como la noción clásica de inmersión regular (EGA IV 16.9.2) resulta inadecuada en el caso no noetheriano, vamos a definir una noción más débil, pero que da lugar a las propiedades deseadas y que equivale a la definición usual en el caso noetheriano.

\textbf{Definición 1.1.} Sean $A$ un anillo conmutativo, $E$ un $A$-módulo proyectivo de tipo finito, y $u:E\longrightarrow A$ un homomorfismo $A$-lineal. Se dice que $u$ es regular si el complejo de Koszul $K_\bullet(u)$ definido por $u$ es acíclico en dimensión $>0$, y por tanto es una resolución de $B=A/u(E)$.

Recordemos que el complejo de Koszul asociado a $u$ se define como sigue: el $A$-módulo graduado $K_\bullet(u)$ es el álgebra exterior $\Lambda(E)$, y el operador borde es la multiplicación interior por $u$ considerado como elemento del dual $\check E$: es una antiderivación de grado $-1$ de $\Lambda(E)$, igual a $u$ sobre $E$.

\textbf{Lema 1.2.} Sea $u:E\longrightarrow A$ un homomorfismo regular. Sean $A'$ una $A$-álgebra y $u':E'\longrightarrow A'$ el homomorfismo deducido de $u$ por cambio de base; se supone que $u'$ es regular. Se pone $J=u(E)$, $J'=JA'=u(E')$. Entonces se tienen las relaciones siguientes para todo entero $n\geq0$:
\begin{align}
 J^n/J^{n+1}\otimes_AA'&\xrightarrow{\sim}J'^n/J'^{n+1}, \tag{1.2.1}\\
 \Tor_i^A(J^n/J^{n+1},A')&=0\quad\text{para }i\geq1, \tag{1.2.2}\\
 J^n\otimes_AA'&\xrightarrow{\sim}J'^n, \tag{1.2.3}\\
 \Tor_i^A(J^n,A')&=0\quad\text{para }i\geq1. \tag{1.2.4}
\end{align}

\emph{Demostración.} a) Supongamos primero que los $J^n/J^{n+1}$ son proyectivos sobre $B=A/J$; como $J$ es de tipo finito, son localmente libres de tipo finito sobre $B$. Puesto que las relaciones que hay que demostrar se verifican puntualmente sobre $\Spec(A)$, quedamos reducidos al caso en que $A$ es local, con $u(E)\subset\rad(A)$. Los $J^n/J^{n+1}$ son entonces libres sobre $B$, y las relaciones $\Tor_i^A(J^n/J^{n+1},A')=0$ para $i\geq1$ equivalen a las relaciones $\Tor_i^A(B,A')=0$ para $i\geq1$. Ahora bien, se pueden calcular los $\Tor_i^A$ utilizando la resolución proyectiva $K_\bullet(u)$ de $B$; son los grupos de homología del complejo $K_\bullet(u)\otimes_AA'=K_\bullet(u')$. Como se supone que $u'$ es regular, este último es acíclico en grados $>0$, de donde resultan las relaciones $\Tor_i^A(J^n/J^{n+1},A')=0$ para $i\geq1$.

Se deducen las relaciones $\Tor_i^A(A/J^n,A')=0$ para $i\geq1$, procediendo por inducción sobre $n$ gracias a la sucesión exacta de los Tor asociada a la sucesión exacta
\[
 0\longrightarrow J^n/J^{n+1}\longrightarrow A/J^{n+1}\longrightarrow A/J^n\longrightarrow0.
\]
Como para todo $k$ se tiene $J^kA'=J'^k$, de esta sucesión exacta y de $\Tor_1^A(A/J^n,A')=0$ se obtienen las relaciones $J^n/J^{n+1}\otimes_AA'\xrightarrow{\sim}J'^n/J'^{n+1}$. De la sucesión exacta
\[
 0\longrightarrow J^n\longrightarrow A\longrightarrow A/J^n\longrightarrow0,
\]
y de $\Tor_1^A(A/J^n,A')=0$, se deducen las relaciones $J^n\otimes_AA'\xrightarrow{\sim}J'^n$. Finalmente, de la sucesión exacta de los Tor asociada a la sucesión exacta
\[
 0\longrightarrow J^{n+1}\longrightarrow J^n\longrightarrow J^n/J^{n+1}\longrightarrow0,
\]
se deducen las relaciones $\Tor_i^A(J^n,A')=0$ para $i\geq1$, por inducción sobre $n$, siendo evidente el caso $n=0$.

b) Vamos a demostrar ahora que los $J^n/J^{n+1}$ son proyectivos sobre $B=A/J$. Como la cuestión es local sobre $\Spec(A)$, podemos suponer que $E$ es libre, sea $E=(A)^d$, de modo que $u$ está definido por una sucesión $f_1,\ldots,f_d$ de elementos de $A$. Pongamos $A_0=\mathbb Z[T_1,\ldots,T_d]$ y sea $A_0\longrightarrow A$ el homomorfismo que envía $T_i$ a $f_i$; denotemos por $u_0:E_0=A_0^d\longrightarrow A_0$ el homomorfismo definido por la sucesión $(T_1,\ldots,T_d)$, de modo que $u$ se deduce de $u_0$ por cambio de base de $A_0$ a $A$. Como la sucesión $(T_1,\ldots,T_d)$ es regular (en el sentido clásico), el homomorfismo $u_0$ es regular, y estamos bajo las condiciones del lema, aplicado al cambio de base $A_0\longrightarrow A$. Ahora bien, los $J_0^n/J_0^{n+1}$ son libres de tipo finito sobre $\mathbb Z$; por tanto, se pueden aplicar los resultados de a). En consecuencia, $J^n/J^{n+1}\xleftarrow{\sim}J_0^n/J_0^{n+1}\otimes_{A_0}A$, y $J^n/J^{n+1}$ es libre de tipo finito.

\textbf{Proposición 1.3.} Sean $A$ un anillo, $E$ un $A$-módulo proyectivo de tipo finito, $u:E\longrightarrow A$ un homomorfismo regular, y $J=u(E)$. Entonces el ideal $J$ es cuasirregular (EGA IV 16.9.7), es decir (EGA IV 16.9.8):
\begin{enumerate}[label=\roman*)]
\item $J$ es de tipo finito.
\item $J/J^2$ es un $A/J$-módulo proyectivo de tipo finito;
\item el homomorfismo canónico
\end{enumerate}
\begin{equation}\tag{1.3.1}
 \Sym_{A/J}(J/J^2)\longrightarrow \gr_J(A)=\bigoplus_{n\geq0}J^n/J^{n+1},
\end{equation}
es un isomorfismo.

\emph{Demostración.} En la demostración de 1.2 b) se ha visto que $J/J^2$ es proyectivo de tipo finito. Retomemos las notaciones de 1.2 b) (como la cuestión es local sobre $\Spec(A)$, se supone todavía que $E$ es libre de tipo finito). Entonces el homomorfismo canónico (1.3.1) se deduce del homomorfismo análogo para $J_0$ mediante el cambio de base $A_0\longrightarrow A$, por las relaciones 1.2. Ahora bien, este último es un isomorfismo, y por tanto lo mismo vale para (1.3.1).

\textbf{Definición 1.4.} Sean $(X,A)$ un sitio conmutativamente localmente anillado, $E$ un $A$-módulo localmente libre de tipo finito, y $u:E\longrightarrow A$ un homomorfismo $A$-lineal. Se dice que $u$ es regular si el complejo de Koszul $K_\bullet(u)$ definido por $u$ es acíclico en grados $>0$. Sea $J$ un ideal de $A$; se dice que $J$ es regular si, localmente sobre $X$, existen un $A$-módulo libre de tipo finito $E$ y un homomorfismo regular sobreyectivo $E\longrightarrow J$.

Sea $i:Y\longrightarrow X$ una inmersión de preesquemas, y sea $U$ un abierto de $X$ tal que $i(Y)\subset U$ y que $i$ sea una inmersión cerrada de $Y$ en $U$. Se dice que $i$ es regular si el ideal de $\cO_U$ que define $i(Y)$ es regular, condición que no depende de $U$. Si $y\in Y$, se dice que $i$ es regular en $y$ si existe un entorno $V$ de $i(y)$ tal que la inmersión $i^{-1}(V)\longrightarrow V$ sea regular. Para que una inmersión sea regular, es necesario y suficiente, pues, que lo sea en un entorno de cada punto.

Esta definición está en conflicto con EGA IV 16.9.2 y sustituye a esta última. Sin embargo, en el caso noetheriano las dos definiciones son equivalentes (EGA IV 19.5.1). Una inmersión que localmente puede definirse por una sucesión regular es regular; una inmersión regular es cuasirregular (1.3).

\textbf{1.4.1.} Sean $X$ un esquema, $J$ un ideal de tipo finito de $\cO_X$, y $x$ un punto de $\Supp(\cO_X/J)$. Diremos que una familia de secciones $f_1,\ldots,f_d$ de $J$ sobre un entorno de $x$ forma un sistema mínimo de generadores de $J$ en un entorno de $x$ si los $(f_i)_x$ son una familia de generadores de $J_x$ de cardinal mínimo, es decir, si sus imágenes en $J_x\otimes k(x)$ forman una base de este espacio vectorial sobre $k(x)$. Obsérvese que los $f_i$ engendran $J$ en un entorno de $x$.

\textbf{1.4.2.} Decir que una inmersión es regular equivale también a decir que está definida en cierto abierto de $X$ por un ideal de tipo finito $J$, y que, para todo punto $x\in\Supp(\cO_X/J)$ y todo sistema mínimo de generadores de $J$ en un entorno de $x$, el complejo de Koszul asociado a esta sucesión es acíclico en un entorno de $x$. En efecto, sea $f_1,\ldots,f_d$ un sistema mínimo de generadores de $J$ en un entorno de $x$ tal que el complejo de Koszul asociado sea acíclico (existe si se supone regular la inmersión); sea $g_1,\ldots,g_d$ otro sistema mínimo de generadores de $J$. Existe un entorno de $x$ en el que $J$ puede ser engendrado tanto por los $f_i$ como por los $g_j$, y por tanto existe en ese entorno una matriz con coeficientes en $\cO_X$ que permite expresar los $g_j$ en función de los $f_i$. Ahora bien, las imágenes de los $f_i$ (resp. de los $g_j$) en $J/J^2$ forman una base en el punto $x$, y por ello en un entorno de $x$; en consecuencia, esta matriz es invertible en un entorno de $x$. Permite, pues, definir un isomorfismo del complejo de Koszul definido por los $g_j$ sobre el complejo de Koszul definido por los $f_i$, lo que muestra que el complejo de Koszul definido por los $g_j$ es acíclico.

\textbf{1.4.3.} Suponiendo solamente que $J$ es de tipo finito y $J/J^2$ localmente libre de tipo finito, observemos que todo sistema mínimo de generadores de $J$ en un entorno de un punto $y$ de $Y$ da una base de $J/J^2$ en un entorno de $y$. Recíprocamente, si se tiene en un entorno de $y$ una base de $J/J^2$, se puede levantar en un entorno de $y$ a una familia de secciones de $J$, que engendran $J_y$ por Nakayama, y por tanto $J$ en un entorno de $y$ porque $J$ es de tipo finito; además, estos generadores son de cardinal mínimo, puesto que dan una base de $J/J^2$.

\textbf{Proposición 1.5.} (Compárese EGA IV 19.1.5 (ii)). Sea $i:Y\longrightarrow X$ una inmersión. Si $i$ es regular, lo mismo vale para toda inmersión deducida de $i$ por cambio de base plano. Para que $i$ sea regular, es necesario y suficiente que lo sea localmente para la topología fielmente plana cuasicompacta.

\emph{Demostración.} La estabilidad por cambio de base plano es evidente.

Recíprocamente, sea $i'$ la inmersión deducida de $i$ por un cambio de base fielmente plano cuasicompacto. Si $i'$ es regular, entonces el ideal $J$ que define $i$ es de tipo finito, y $J/J^2$ es localmente libre de tipo finito, pues estas propiedades descienden por morfismos fielmente planos cuasicompactos.

Si $f_1,\ldots,f_d$ es una sucesión de secciones de $J$ que levanta una base de $J/J^2$, el complejo de Koszul asociado a las imágenes inversas de los $f_i$ es acíclico en grados $>0$ por 1.4; como es la imagen inversa del complejo de Koszul asociado a los $f_i$, este último es acíclico en grados $>0$, y $i$ es regular.

\textbf{Definición 1.6.} Sea $i:Y\longrightarrow X$ una inmersión definida en un abierto $U$ de $X$ por un ideal $J$ de $\cO_X$; el $\cO_Y$-módulo $J/J^2$ (independiente de $U$) se llama haz conormal de $Y$ en $X$, y se denota por $N_{X/Y}$. Si $i$ es regular, $J/J^2$ es localmente libre de tipo finito, y su rango se llama codimensión de $Y$ en $X$, o también codimensión de $i$.

Si se tienen dos inmersiones cerradas $Z\xrightarrow{j}Y\xrightarrow{i}X$, los haces conormales forman una sucesión exacta
\begin{equation}\tag{1.6.1}
 j^*(N_{X/Y})\longrightarrow N_{X/Z}\longrightarrow N_{Y/Z}\longrightarrow0.
\end{equation}
Sean, en efecto, $I$ el ideal de $Y$ en $X$, $J$ el ideal de $Z$ en $X$, y $K$ el ideal de $Z$ en $Y$. Se tiene la sucesión exacta
\[
0\longrightarrow I\longrightarrow J\longrightarrow K\longrightarrow0;
\]
de donde, tensorizando por $\cO_X/I=\cO_Y$, resulta la sucesión exacta
\[
 I/I^2\longrightarrow J\otimes_{\cO_X}\cO_Y\longrightarrow K\otimes_{\cO_X}\cO_Y\longrightarrow0;
\]
tensorizando por $\cO_Z=\cO_X/J=\cO_Y/K$, se obtiene la sucesión exacta
\[
 j^*(I/I^2)\longrightarrow J\otimes_{\cO_X}\cO_Z\longrightarrow K\otimes_{\cO_Y}\cO_Z\longrightarrow0,
\]
que es la sucesión buscada.

\textbf{Proposición 1.7.} (Compárese EGA IV 19.1.5 (iv)). Sean $i:Y\longrightarrow X$ y $j:Z\longrightarrow Y$ dos inmersiones:
\begin{enumerate}[label=\roman*)]
\item Si $i$ y $j$ son regulares, $i\circ j$ es regular, y la sucesión
\[
0\longrightarrow j^*(N_{X/Y})\longrightarrow N_{X/Z}\longrightarrow N_{Y/Z}\longrightarrow0
\]
es exacta; en particular, $\codim(i\circ j)=\codim(i)+\codim(j)$.
\item Si $i\circ j$ e $i$ son regulares, y si la sucesión
\[
0\longrightarrow j^*(N_{X/Y})\longrightarrow N_{X/Z}\longrightarrow N_{Y/Z}\longrightarrow0,
\]
es exacta y localmente escindida, $j$ es regular;
\item Si $X$ es noetheriano, y si $i\circ j$ y $j$ son regulares, $i$ es regular en los puntos de $Z$.
\end{enumerate}

\emph{Demostración.} Como la cuestión es local sobre $X$, podemos suponer que $X$ es afín y que las inmersiones son cerradas; sean $X=\Spec(A)$, $Y=\Spec(B)$, $Z=\Spec(C)$, con $B=A/I$, $C=A/J=B/K$.

i) Localizando si es necesario, podemos suponer que $I/I^2$ es libre sobre $B$ y $K/K^2$ libre sobre $C$ (1.3). Sean $f_1,\ldots,f_n$ elementos de $A$ que levantan una base de $I/I^2$, y $g_1,\ldots,g_m$ elementos de $A$ que levantan una base de $K/K^2$; localizando más si es necesario, podemos suponer que $f_1,\ldots,f_n$ engendran $I$, que $f_1,\ldots,f_n,g_1,\ldots,g_m$ engendran $J$, y que $\bar g_1,\ldots,\bar g_m$ engendran $K$, siendo $\bar g_j$ la imagen de $g_j$ en $B$ (1.4); además, $f_1,\ldots,f_n$ forman un sistema mínimo de generadores de $I$, y $\bar g_1,\ldots,\bar g_m$ un sistema mínimo de generadores de $K$.

Consideremos el diagrama
\begin{equation}\tag{1.7.1}
\begin{tikzcd}[column sep=large,row sep=large]
A_0=\mathbb Z[T_1,\ldots,T_{n+m}] \arrow[r] \arrow[d] & A \arrow[d]\\
B_0=\mathbb Z[T_1,\ldots,T_m] \arrow[r] \arrow[d] & B \arrow[d]\\
C_0=\mathbb Z \arrow[r] & C,
\end{tikzcd}
\end{equation}
donde las flechas horizontales se definen enviando $T_i$ a $f_i$ para $1\leq i\leq n$, y $T_{n+j}$ a $g_j$ para $1\leq j\leq m$; el homomorfismo $A_0\longrightarrow B_0$ se define enviando $T_i$ a $0$ para $1\leq i\leq n$, y $T_{n+j}$ a $T_j$ para $1\leq j\leq m$; y el homomorfismo $B_0\longrightarrow C_0$ anulando los $T_j$, $1\leq j\leq m$. Se tiene entonces
\[
 B=A\otimes_{A_0}B_0,
 \qquad
 C=A\otimes_{A_0}C_0=B\otimes_{B_0}C_0.
\]
La situación estudiada se deduce, pues, de la situación $(A_0,B_0,C_0)$ mediante el cambio de base $A_0\longrightarrow A$. Ahora bien, las inmersiones correspondientes $i_0$, $j_0$ e $i_0\circ j_0$ son regulares; por tanto, los grupos de homología de los complejos de Koszul asociados a las sucesiones $(f_1,\ldots,f_n)$, $(f_1,\ldots,f_n,g_1,\ldots,g_m)$ y $(\bar g_1,\ldots,\bar g_m)$ son respectivamente los $\Tor_i^{A_0}(B_0,A)$, $\Tor_i^{A_0}(C_0,A)$, $\Tor_i^{A_0}(C_0,B)$. Por consiguiente:
\begin{equation}\tag{1.7.2}
\begin{aligned}
 i\ \text{regular} &\Longleftrightarrow \forall i\geq1,
 \Tor_i^{A_0}(B_0,A)=0,\\
 i\circ j\ \text{regular} &\Longleftrightarrow \forall i\geq1,
 \Tor_i^{A_0}(C_0,A)=0,\\
 j\ \text{regular} &\Longleftrightarrow \forall i\geq1,
 \Tor_i^{B_0}(C_0,B)=0.
\end{aligned}
\end{equation}
Utilizamos entonces la sucesión espectral de transitividad ([1], XVI 5)
\begin{equation}\tag{1.7.3}
E^2_{pq}=\Tor_p^{B_0}\bigl(C_0,\Tor_q^{A_0}(B_0,A)\bigr)
\Longrightarrow \Tor_n^{A_0}(C_0,A).
\end{equation}
Por las hipótesis de i), $E^2_{pq}=0$ para $p$ o $q\ne0$. Por tanto, degenera, y $i\circ j$ es regular por (1.7.2).

Por (1.6.1), la sucesión
\[
 j^*(N_{X/Y})\longrightarrow N_{X/Z}\longrightarrow N_{Y/Z}\longrightarrow0
\]
es exacta. Como $i\circ j$ es regular, los tres haces conormales son localmente libres de tipo finito; además, su rango en un punto es el número de elementos de un sistema mínimo de generadores del ideal correspondiente; por la demostración precedente, se tiene
\[
 \rang(N_{X/Z})=\rang(N_{X/Y})+\rang(N_{Y/Z}).
\]
Se sigue que el núcleo de $j^*(N_{X/Y})\longrightarrow N_{X/Z}$ es localmente libre de rango cero, y por tanto que la sucesión (\sisi{1.8.1}{1.6.1}) es exacta.

ii) Podemos suponer que $J/J^2$ es libre sobre $C$, e $I/I^2$ libre sobre $B$; además, la hipótesis hecha sobre la sucesión de los haces conormales permite suponer que $I/I^2\otimes_A A/J$ es factor directo de $J/J^2$. Podemos encontrar, pues, elementos $f_1,\ldots,f_n$ de $A$ que levantan una base de $I/I^2$, y elementos $g_1,\ldots,g_m$ de $A$ tales que $f_1,\ldots,f_n,g_1,\ldots,g_m$ levantan una base de $J/J^2$, y cuyas imágenes $g_1,\ldots,g_m$ engendran $K/K^2$; localizando de nuevo si es necesario, podemos suponer que $f_1,\ldots,f_n$ forman un sistema mínimo de generadores de $I$, que $f_1,\ldots,f_n,g_1,\ldots,g_m$ forman un sistema mínimo de generadores de $J$, y que $\bar g_1,\ldots,\bar g_m$ engendran $K$. Podemos entonces construir el diagrama (1.7.1) y utilizar la sucesión espectral (1.7.3).

Por (1.7.2), como $i$ es regular, se tiene $E^2_{pq}=0$ para $q\geq1$; se obtienen, pues, isomorfismos de borde $E^2_{p,0}\simeq\Tor_p^{A_0}(C_0,A)$. Como $i\circ j$ es regular, esto implica por (1.7.2) que $E^2_{p,0}=0$ para $p\geq1$, y por tanto que $j$ es regular, siempre por (1.7.2).

iii) Esta última propiedad está demostrada en EGA IV 19.1.5 (iv), puesto que en el caso noetheriano la definición de EGA IV coincide con la que adoptamos aquí.

\textbf{Proposición 1.8.} Sean \(X\) un esquema, \(E\) un \(\cO_X\)-módulo localmente libre de tipo finito, \(u:E\to\cO_X\) un homomorfismo regular, y \(J=u(E)\). Denotamos por \(Y\) el subesquema cerrado de \(X\) definido por \(J\), por \(X'\) el esquema obtenido explotando \(Y\), por \(P=\bfP(E)\) el fibrado proyectivo asociado a \(E\), y por \(p:P\to X\) el morfismo estructural. Sobre \(P\), consideramos la sucesión exacta canónica
\begin{equation}\tag{1.8.1}
0\to F\to p^*(E)\to\cO_P(1)\to0,
\end{equation}
y el morfismo compuesto
\[
 v:F\longrightarrow p^*(E)\xrightarrow{p^*(u)}\cO_P .
\]
Entonces:
\begin{enumerate}[label=\roman*)]
\item La formación de \(X'\) conmuta con todo cambio de base \(f:T\to X\) tal que \(f^*(u)\) sea regular.
\item La inmersión cerrada \(j:X'\to P\) definida por el homomorfismo sobreyectivo \(\cS(E)\longrightarrow\bigoplus_{n\ge0}J^n\), definido por \(u\) en grado \(1\), es una inmersión regular, siendo su ideal la imagen del morfismo regular \(v\); en particular, si \(\operatorname{codim}(i)\) es constante, \(\operatorname{codim}(j)=\operatorname{codim}(i)-1\).
\item El \(\cS(E)\)-módulo \(\bigoplus_{n\ge0}J^n\) es perfecto como \(\cS(E)\)-módulo graduado, es decir, existe localmente una resolución finita y graduada de \(\bigoplus_{n\ge0}J^n\) por \(\cS(E)\)-módulos graduados localmente libres de tipo finito (VI 1.2).
\end{enumerate}

i) Como \(X'=\Proj(\bigoplus_{n\ge0}J^n)\), el resultado es una consecuencia inmediata de las relaciones (1.2.3) y de EGA II.

ii) Como las propiedades enunciadas son locales sobre \(X\), podemos suponer que \(X\) es afín y, además, que \(E\) es libre sobre \(X\). Ponemos \(X=\Spec(A)\). Hacemos primero la demostración en el caso estándar \(A=\mathbb Z[T_1,\ldots,T_d]\), estando \(E\) provisto de una base \(e_1,\ldots,e_d\), con \(u(e_i)=T_i\). Entonces \(P=\Proj\bigl(\mathbb Z[T_1,\ldots,T_d,X_1,\ldots,X_d]\bigr)\), tomando la graduación respecto de los \(X_i\). Para demostrar que \(v\) es un morfismo regular cuya imagen es el ideal de \(j\), basta hacerlo localmente sobre \(P\), y por tanto sobre los abiertos \(D_+(X_i)\); situémonos, por ejemplo, sobre \(D_+(X_1)\). Entonces \(p^*(E)\) está definido por el módulo
\[
 \mathbb Z[T_1,\ldots,T_d,X_2/X_1,\ldots,X_d/X_1]\otimes E,
\]
y \(\cO_P(1)\) por
\[
 \mathbb Z[T_1,\ldots,T_d,X_2/X_1,\ldots,X_d/X_1],
\]
enviando el morfismo canónico \(e_i\) a \(X_i/X_1\) para todo \(i\). Entonces \(F\) admite como base los elementos \(e_i-e_1.X_i/X_1\) para \(i>1\).
Se tiene \(v(e_i-e_1.X_i/X_1)=T_i-T_1.X_i/X_1\); pongamos \(T'_i=T_i-T_1X_i/X_1\).
Como la sucesión de los \(T_i\) es regular en \(\mathbb Z[T_1,\ldots,T_d,X_2/X_1,\ldots,X_d/X_1]\), lo mismo vale para la sucesión \(T_1,T'_2,\ldots,T'_d\), que se deduce de ella mediante una matriz invertible, y por tanto también para la sucesión \(T'_2,\ldots,T'_d\) extraída de ella (estamos en un caso noetheriano). Por consiguiente, \(v\) es un morfismo regular.

Hay que demostrar, pues, que los \(T'_i\) engendran el ideal de \(j\). En \(\bigoplus_{n\ge0}J^n\), denotaremos, para \(m\le p\), por \(x^{(m)}\) el elemento \(x\) de \(J^p\), considerado de grado \(m\). Sobre \(D_+(X_1)\), \(X'=\Spec\left(\left(\bigoplus_{n\ge0}J^n\right)_{(T_1^{(1)})}\right)\), y \(j\) está definido por el morfismo que envía \(X_i\) a \(T_i^{(1)}\). En este morfismo, la imagen de \(T'_i\) es \(T_i^{(0)}-T_1^{(0)}.T_i^{(1)}/T_1^{(1)}=\bigl((T_iT_1)^{(1)}-(T_1T_i)^{(1)}\bigr)/T_1^{(1)}=0\). Por tanto, los \(T'_i\) pertenecen al ideal de \(j\). Como todo elemento de \(\sisi{\mathbb Z[T_1,\ldots,T_d,X_2/X_1,\ldots,X_d,X_1]}{\mathbb Z[T_1,\ldots,T_d,X_2/X_1,\ldots,X_d/X_1]}\) puede expresarse como polinomio respecto de \(T_1,T'_2,\ldots,T'_d,X_2/X_1,\ldots,X_d/X_1\), para demostrar que los \(T'_i\) engendran el ideal de \(j\) basta demostrar que ningún elemento no nulo de \(\mathbb Z[T_1,X_2/X_1,\ldots,X_d/X_1]\) pertenece al ideal de \(j\), lo que equivale a demostrar que el homomorfismo
\[
\mathbb Z[T_1,X_1,\ldots,X_d]\longrightarrow\bigoplus_{n\ge0}J^n
\]
es inyectivo, pues \(T_1^{(1)}\) no es divisor de 0 en \(\bigoplus_{n\ge0}J^n\). Si se gradúa \(\mathbb Z[T_1,X_1,\ldots,X_d]\) asignando el grado \(0\) a \(T_1\), y el grado \(1\) a los \(X_i\), el núcleo del homomorfismo considerado es un ideal graduado. Ahora bien, está claro que este núcleo está contenido en el del homomorfismo compuesto
\[
\mathbb Z[T_1,X_1,\ldots,X_d]\longrightarrow
\bigoplus_{n\ge0}J^n\longrightarrow A,
\]
donde el último está definido en grado \(n\) por la inyección de \(J^n\) en \(A\). Ahora bien, el núcleo del homomorfismo compuesto es el ideal principal engendrado por \((X_1-T_1)\), que por razones de integridad no contiene elementos homogéneos no nulos. Por tanto, el primer homomorfismo es inyectivo, lo que termina la demostración en el caso considerado.

Volvamos al caso general, donde \(A\) es arbitrario; todavía podemos suponer que \(E\) es libre sobre \(X\). Afectemos con el índice \(0\) el caso particular recién tratado; entonces, por la parte i) de la proposición, la situación considerada se obtiene mediante el cambio de base \(A_0\to A\) a partir de la situación estudiada anteriormente, estando definido el homomorfismo \(A_0\to A\) al tomar una base de \(E\), su imagen \(f_1,\ldots,f_d\) en \(J\), y enviar \(T_i\) a \(f_i\). Pongamos \(B_0=\left(\bigoplus_{n\ge0}J_0^n\right)_{(T_1^{(1)})}\), \(B=\left(\bigoplus_{n\ge0}J^n\right)_{(f_1^{(1)})}\); entonces \(B=B_0\otimes_{A_0}A\). El complejo de Koszul definido por \(v:F\to\cO_P\) se deduce por cambio de base del complejo de Koszul definido por \(v_0\); como este último es una resolución proyectiva de \(B_0\) (situándose sobre un abierto conveniente), los grupos de homología de \(K_\bullet(v)\) son, sobre este abierto, los \(\Tor_i^{A_0}(A,B_0)\). Estos últimos son las componentes de grado \(0\) de los \(\Tor_i^{A_0}\left(A,\left(\bigoplus_{n\ge0}J_0^n\right)_{T_1^{(1)}}\right)=\left(\bigoplus_{n\ge0}\Tor_i^{A_0}(A,J_0^n)\right)_{T_1^{(1)}}\). Ahora bien, los \(\Tor_i^{A_0}(A,J_0^n)\) son nulos para \(i\ge1\) por las relaciones 1.2; por tanto, \(K_\bullet(v)\) es una resolución de \(B\), lo que demuestra que \(j\) es regular.

iii) Como la propiedad es local sobre \(X\), todavía podemos suponer que \(X\) es afín y \(E\) libre de tipo finito. Retomamos entonces las notaciones de ii). Consideremos primero el caso estándar de ii). El anillo \(A_0\) es entonces noetheriano y de dimensión cohomológica finita, y el álgebra \(\cS(E_0)\) también. Se puede construir paso a paso una resolución parcial graduada de \(\bigoplus_{n\ge0}J_0^n\), por \(\cS(E_0)\)-módulos graduados libres de tipo finito, de longitud superior a la dimensión cohomológica de \(A_0\); el último núcleo es entonces un \(\cS(E_0)\)-módulo graduado, proyectivo de tipo finito como \(\cS(E_0)\)-módulo no graduado, y por tanto localmente libre de tipo finito sobre \(\Spec(A_0)\) como \(\cS(E_0)\)-módulo graduado, por VI 1.5.1. Se obtiene entonces la resolución finita anunciada sustituyendo el último módulo de la resolución parcial por el núcleo correspondiente.

Si se tensoriza por \(A\) la resolución así construida, se obtiene un complejo graduado de \(\cS(E)\)-módulos graduados localmente libres de tipo finito, y hay que ver que es una resolución de \(\bigoplus_{n\ge0}J^n\). Ahora bien, el \(i\)-ésimo módulo de homología del complejo así obtenido es un \(\cS(E)\)-módulo graduado cuya componente \(n\)-ésima es \(\Tor_i^{A_0}(A,J_0^n)\), que es nula para \(i\ge1\) por (1.2.4). Por consiguiente, por (1.2.3), se ha construido una resolución del tipo deseado.

\textbf{Proposición 1.9.} Sean \(i:Y\to X\) una inmersión cerrada regular, \(X'\) el \(X\)-esquema obtenido explotando \(Y\), y \(f:X'\to X\) el morfismo estructural. Entonces \(i\) y \(f\) son morfismos perfectos (III 4.1).

Como la propiedad es local sobre \(X\), se pueden suponer dados un \(\cO_X\)-módulo localmente libre de tipo finito \(E\), y un homomorfismo regular \(u:E\to\cO_X\), cuya imagen es el ideal \(J\) de \(Y\). El complejo de Koszul \(K_\bullet(u)\) es entonces una resolución finita de \(\cO_Y\) por \(\cO_X\)-módulos localmente libres de tipo finito; por consiguiente, \(\cO_Y\) es un \(\cO_X\)-módulo perfecto. Se sigue que toda inmersión regular es un morfismo perfecto. Por 1.8, \(f\) puede escribirse como el compuesto de una inmersión regular \(X'\to\bfP(E)\) y del morfismo estructural \(\bfP(E)\to X\), que es liso, y por tanto perfecto; es, pues, el compuesto de dos morfismos perfectos y, en consecuencia, un morfismo perfecto.

\textbf{Proposición 1.10.} Sean \(X\), \(Y\) dos esquemas lisos sobre \(S\). Entonces toda \(S\)-inmersión de \(X\) en \(Y\) es regular. En particular, toda sección de un morfismo liso es una inmersión regular.

La proposición resulta inmediatamente de EGA IV 17.12.1, puesto que una inmersión regular en el sentido de EGA IV es regular en el sentido de 1.4.
